% Physical page 2 of the concise study, which carries this sheet and nothing
% else. Every row is the generated orientation-v1 projection of
% research/chronology.toml, placed once in the generator's order, exactly as
% the reviewed Scriptural Date and Location appendix of the expansive study
% places it; the place phrases are that appendix's. The generated definitions
% are imported here, once, because the concise entrypoint's own preamble does
% not carry them.
\input{liturgy/roman-rite/postconciliar/roman-missal-third-edition-en-us-2011/propers/temporal/pc-s55-twenty-ninth-sunday-in-ordinary-time-year-a/research/chronology-annotations}

\section{Scriptural Date and Location}

{\small\noindent Each passage the Mass reads or sings from Scripture, in
canonical order: what is dated, where it belongs, and the corpus date. Rows
sharing one answer name every element that uses it. The antiphons' psalms are
cited by Vulgate number (Pss 16, 32), the responsorial psalm by the
Lectionary's (Ps 96).\par}
\medskip

\zlabel{triptych:concise:chronology:start}
\begin{chronologyorientationtable}
\chronologyorientationrow{row-29cbe0d234a47aa8}{Occasion and place unknown; the psalms' titles name none.}
\addlinespace
\chronologyorientationrow{row-6e9c8175463dbd35}{Place of composition unknown.}
\addlinespace
\chronologyorientationrow{row-d7a88409f6dff8dd}{Ascribed to David by its Latin and Greek title; David reigns at Hebron, then at Jerusalem. The reign dates David, not the psalm.}
\addlinespace
\chronologyorientationrow{row-b6be5603e767096a}{The house rebuilt after the captivity, named in the psalm's Latin and Greek title: the Temple at Jerusalem.}
\addlinespace
\chronologyorientationrow{row-7a3e8d424af7908c}{The Messiah's coming that the Fathers read in this psalm, taken as the Nativity at Bethlehem.}
\midrule
\chronologyorientationrow{row-2da3a3c2bfe61ef2}{Isaias, a citizen of Jerusalem, prophesies concerning Juda and Jerusalem.}
\addlinespace
\chronologyorientationrow{row-e535d564c1cee995}{An unnamed prophet, place not stated; his oracles bid the exiles go out of Babylon.}
\midrule
\chronologyorientationrow{row-a180ccaa931e671d}{Jesus answers the tribute question while teaching in the temple at Jerusalem.}
\addlinespace
\chronologyorientationrow{row-5396102e03d2c0da}{Matthew writes for Jewish readers; the place is not fixed (Eusebius: on leaving Palestine).}
\midrule
\chronologyorientationrow{row-4b690f3afb3478b9}{Mark writes at Rome for Roman Christians (Chrysostom names Egypt).}
\midrule
\chronologyorientationrow{row-11f6d60f0fa3c25e}{Paul writes from captivity at Rome (some say Caesarea) to the church at Philippi.}
\midrule
\chronologyorientationrow{row-d67f41bffc530c4a}{Paul writes from Corinth (an ancient subscription says Athens) to the church at Thessalonica.}
\end{chronologyorientationtable}
\zlabel{triptych:concise:chronology:end}

\medskip
{\footnotesize\noindent The Isaiah rows date the prophet's ministry in one
traditional account and the critical placement of chapters 40--55, not the
oracle; the composition rows date the books, and the Psalter's bound dates no
psalm singly. Sources are given passage by passage in the full study's appendix
Scriptural Date and Location.\par}
